%======================================================================
\section{引言}
\label{sec:intro}
%======================================================================

由大语言模型（LLM）驱动的多智能体系统日益承担起\emph{集体判断}任务：多个模型实例独立分析问题、交换立场，并聚合为共识决策~\cite{han2024llm,tran2025multiagent,wang2023selfconsistency,du2024improving}。然而，对此类系统的评估至今仍是一维的：只看最终决策的准确率。评审团\emph{如何}抵达决策——它是否保留了正当的异议、维持了评分多样性、是否基于正确的理由收敛——却很少被测量，因为自由文本形式的审议使这些行为量难以提取。

这种忽视是危险的。人类委员会存在羊群效应是众所周知的：评审人之间的讨论与分数压缩和从众效应相关，初始多数派立场会不成比例地决定最终结果。LLM 评审团是否表现出类似的病理——其方向和量级如何——在很大程度上尚未得到刻画。这一问题对机制设计至关重要：若羊群效应确实存在，其结构决定了哪一类干预能够修复它。如果病理是静态偏置，聚合时校正即可；如果讨论本身摧毁信息，则需要行为调控；如果审议只在部分实例上有所回报，那么路由才是合适的方案。

回答上述任何一个问题都需要传统管道所缺乏的仪器：羊群效应必须先变得\emph{可测量}，才能被治疗。本文表明，此前开发的博弈论共识框架恰好提供了这些仪器。其 \textbf{AEIC 分解}（第~\ref{sec:aeic}~节）将每个智能体的输出结构化为假设、证据、推论与一个显式结论字段——一枚立场探针，使集体行为可被直接测量（比非结构化文本的提取可靠性高 28 个百分点）。其\textbf{三层 Stackelberg 激励机制}（第~\ref{sec:game}~节）具有形式化的 $\epsilon$-激励相容保证（第~\ref{sec:theory}~节），恰好映射到三个干预位置：聚合时的静态参数优化、讨论中的在线与动态调控、以及逐实例的\textbf{自适应协议路由}（第~\ref{sec:topology}~节）。在问答基准上，框架已回答了\emph{决策质量}问题（路由器实现 87.6\% 共识准确率）；本文则让它回答\emph{行为健康}问题——以自动化学术评审为试验场，这是一个高风险且本质上关乎行为的领域，LLM 智能体已在其中为研究论文生成评审意见~\cite{liang2024llmreview,jin2024agentreview}。

验证遵循\emph{诊断后干预}设计（第~\ref{sec:exp}~节）。\textbf{研究一}汇集了 300 篇分层采样的 ICLR~2024 论文（明确接收/边界/明确拒稿），由 5 个同质 LLM 评审智能体组成的评审团在三种条件下评审——独立聚合、三轮讨论与元评审综合。诊断以强有力的统计证据确认了三种病理：独立评审中的系统性\emph{拒稿偏置}（明确接收层论文接收率仅 14.4\%）；讨论中的\emph{负向棘轮}——少数派立场以 96:0 向拒稿方向翻转（$p{=}2.5{\times}10^{-29}$），消灭率 78\%，评分方差坍缩 51\%（$p{=}3.9{\times}10^{-32}$）；以及讨论完全没有带来决策收益（一致率 61.7\%$\to$59.7\%，不显著）。在 DeepSeek-V3、GLM-4-Flash 与 DeepSeek-R1 上的跨模型复现表明，羊群效应是普遍的，但其\emph{方向}依赖于模型：V3 与 R1 评审团向拒稿方向棘轮化，而 GLM 评审团向接收方向棘轮化；且推理能力最强的 R1 遭受最大的多样性坍缩（76\%）——能力并不能抵御羊群效应。

\textbf{研究二}随后将框架的三个层次部署为针对性干预，每层对应一种病理。静态层（\emph{软校准聚合}）将决策一致率从 60\% 提升至 67\%（双臂 66.7\%/67.3\%，合并 McNemar $p{=}0.034$），明确接收层从 2\% 修复至 44\%，隐含阈值 4.88 跨臂完全一致；Stackelberg 优化器自主学到恒正校准项（$\tau{=}+1.04$），证明校准是均衡的内生属性。该精度距人类程序委员会的委员会间一致性（约 74\%，NeurIPS 一致性实验）仅 7 个百分点——决议噪声构成固有测量天花板。动态层（\emph{证据锚定的行为调控}）基本未能撼动负向棘轮（消灭率 78\%$\to$71\%；被干预少数派的坚持率 29\% 对 28\%），但保持了决策质量——这一诚实的负结果表明提示层面的调控无法克服由能力驱动的从众。路由层（\emph{逐论文协议选择}）以低 56\% 的推理成本匹配最优固定协议（64.7\%），将审议确立为按需调用的资源而非常规默认。

我们的贡献如下：
\begin{enumerate}
    \item \textbf{LLM 评审团的大规模配对诊断。}以配对设计和强有力的统计证据确立三种病理——拒稿偏置、负向棘轮、多样性坍缩；所采用的行为健康指标（少数派存活、方差保持、漂移）被证明远比决策准确率敏感。
    \item \textbf{具有相反病态吸引子的跨模型普遍性。}羊群效应的方向依赖于模型（拒稿棘轮 vs. 接收谄媚），且更强的推理能力与更强的多样性坍缩共现。
    \item \textbf{校准优先的干预原则。}聚合时校准修复决策偏置（+6~7 个百分点）并逼近人类委员会间一致性定义的测量天花板；提示层面的行为调控无法恢复多样性；逐论文路由在低 56\% 成本下保持准确率——三者共同描绘了哪类机制治疗哪种病理的图谱。
\end{enumerate}
